% ch25-events-notifications-audit
\section{الأحداث الدوميّة، الإشعارات، والتدقيق (\en{Audit}) — وليه بنعمل كده؟}

في الفصل ده هنجاوب على سؤال بسيط في ظاهره، عميق في باطنه: لمّا يحصل «شيء مهمّ» في
النظام — اتأكّد حجز، وصلت رسالة، اتقدّم طالب على منحة — إزاي بنخلّي بقيّة النظام
«يعرف» ويردّ، من غير ما نكدّس كل ردود الفعل دي جوّه المنطق الأساسيّ؟ الجواب اسمه
\en{Domain Event} (الحدث الدوميّ). وهنكمّل بعدها على أخته التوأم في المساءلة:
\en{Audit} (التدقيق) — إزاي بنسجّل «مين عمل إيه وامتى» على كل فعل حسّاس.

القارئ هنا يُفترَض إنه عارف من قبل ثلاث حاجات، فمش هنعيد شرحها بعمق: \en{MediatR}
كناقل رسائل داخليّ (شفناه في الجزء الأوّل)، والإشعارات على مستوى الاستخدام —
يعني إزاي \code{INotificationDispatcher} بيوصّل الإشعار للجرس والبريد (فصل الإشعارات)،
والمعترِض على \code{SaveChanges} اللي بيلفّ عمليّة الحفظ. اللي جديد النهارده هو
الحلقة اللي بتربط دول ببعض: من لحظة وقوع الحدث لحدّ ظهور نتيجته.

\begin{idea}[الفكرة: الحدث يعلن، والمشترِكون يردّون]
بدل ما «معالِج الحجز» يعرف كل حاجة عن الإيميل و\en{SignalR} والتحليلات، هو بس
بيعلن: «حصل حجز مؤكَّد». مين يهتمّ يشترك ويردّ. ده نفس مبدأ «افصل مَن يُعلِن
عن مَن يتصرّف»، وهو اللي بيخلّي كل قطعة صغيرة، مستقلّة، وسهلة الاختبار.
\end{idea}

% ─────────────────────────────────────────────────────────────────────────────
\subsection{يعني إيه «حدث دوميّ» (\en{Domain Event})؟}

الحدث الدوميّ هو مجرّد إعلان: «حاجة ذات معنى في مجال العمل حصلت خلاص». مش أمر
(\en{command}) بيطلب تنفيذ، ولا استعلام (\en{query}) بيطلب بيانات — إعلان بحقيقة
وقعت. في \en{ScholarPath} الأحداث دي \en{records} صغيرة ثابتة (\en{immutable})،
كلّها بترث نوعًا واحدًا اسمه \code{DomainEvent}، واللي بدوره بيحقّق واجهة
\code{INotification} بتاعة \en{MediatR} — يعني \en{MediatR} يعرف ينشره لكل مشترِك:

\begin{english}
\begin{lstlisting}[language={[Sharp]C}]
// every event inherits DomainEvent -> which is a MediatR.INotification
public abstract record DomainEvent : MediatR.INotification;

public record ChatMessageReceivedEvent(
    Guid ConversationId,
    Guid MessageId,
    Guid SenderId,
    Guid RecipientId,
    string BodyPreview) : DomainEvent;

public record BookingConfirmedEvent(
    Guid BookingId, Guid StudentId, Guid ConsultantId) : DomainEvent;

public record ApplicationSubmittedEvent(
    Guid ApplicationId, Guid StudentId, Guid? ScholarshipId) : DomainEvent;
\end{lstlisting}
\end{english}

هنشوف إيه في المقتطف ده:
\begin{itemize}
\item \code{DomainEvent} نوع مجرّد (\en{abstract record}) بيرث \code{INotification}
  — دي التذكرة اللي بتخلّي \en{MediatR} قادر ينشر أيّ حدث لكل مَن اشترك فيه.
\item كل حدث بيحمل الحقائق بس، مش السلوك: \code{ChatMessageReceivedEvent}
  فيه مُعرِّف المحادثة والرسالة والمُرسِل والمستقبِل ومعاينة نصّ — لا إيميل ولا
  دفع ولا أيّ «ردّ فعل».
\item الأسماء دايمًا بصيغة الماضي: \code{BookingConfirmed}،
  \code{ApplicationSubmitted}، \code{ChatMessageReceived} — لأنها بتوصف حاجة
  حصلت فعلًا، مش حاجة عايزين تحصل.
\end{itemize}

\begin{keyp}[نقطة مهمة: حدث بيحمل حقائق، مش سلوك]
لو لقيت نفسك بتحطّ منطق إرسال بريد جوّه تعريف الحدث، فأنت خلطت الإعلان بالتصرّف.
الحدث «رسالة في زجاجة» — بيقول ما حصل وخلاص. التصرّف مكانه المشترِك، وهنشوفه دلوقتي.
\end{keyp}

% ─────────────────────────────────────────────────────────────────────────────
\subsection{بعد الحفظ فقط: ليه ننشر الحدث بعد \code{SaveChanges} الناجح؟}

في قاعدة ذهبيّة: الحدث ما يتنشرش إلّا بعد ما \code{SaveChangesAsync} ينجح
ويثبّت التغيير في قاعدة البيانات. ليه؟ لأنّ ردّ الفعل بيبني على «حقيقة». لو نشرنا
\code{ChatMessageReceivedEvent} قبل الحفظ، وبعدين الحفظ فشل أو اترجّع
(\en{rollback})، هنكون بعتنا إشعار وبريد عن رسالة مش موجودة أصلًا. ده كذب
على المستخدم، وخطأ يصعُب تتبّعه.

عشان كده الترتيب في المعالِج ثابت: احفظ الأوّل، وبعد نجاح الحفظ ابعت الدفع الحيّ،
وبعدها انشر الحدث. المقتطف التالي من \code{SendMessageCommandHandler} بيوضّح
اللحظة دي بالضبط (مقتصّ):

\begin{english}
\begin{lstlisting}[language={[Sharp]C}]
// 1) persist the truth first: the message row + conversation pointers
await db.SaveChangesAsync(ct);

// 2) live push to whoever has the conversation open now (SignalR group)
await chatNotifier.NotifyNewMessageAsync(
    conversation.Id, message.Id, message.SenderId, message.Body, message.SentAt, ct);

// 3) THEN, and only then: publish the domain event (bell + email for the absent party)
await publisher.Publish(
    new ChatMessageReceivedEvent(
        ConversationId: conversation.Id,
        MessageId:      message.Id,
        SenderId:       senderId,
        RecipientId:    request.RecipientId,
        BodyPreview:    sanitizedBody),
    ct);
\end{lstlisting}
\end{english}

هنشوف إيه:
\begin{itemize}
\item السطر \code{publisher.Publish(new ChatMessageReceivedEvent(...))} هو النقطة
  اللي بيتفرّع منها كل شيء. \code{publisher} هو \code{IPublisher} بتاع
  \en{MediatR}؛ بمجرّد ما تناديه، \en{MediatR} بيروح يدوّر على كل مشترِك في
  \code{ChatMessageReceivedEvent} ويناديه.
\item لاحظ إنّ الحفظ (\code{SaveChangesAsync}) قبله بسطور. لو رمى استثناء، السطر
  ده أصلًا ما بيتنفّذش — فمفيش إشعار عن رسالة اترجّعت.
\end{itemize}

\begin{why}
ليه ننشر بعد الحفظ مش قبله؟ عشان «الاتّساق»: كل ردّ فعل (إشعار، بريد، تحليلات)
لازم يعكس حالة حقيقيّة ثابتة في القاعدة. النشر بعد النجاح بيضمن إنّ أيّ ردّ
فعل يشوف نفس الحقيقة اللي شافها العميل. لو عكسنا الترتيب، هنعيش في احتمال دائم إننا
نبلّغ عن حاجة اتلغت — وده أسوأ من إشعار متأخّر بثوانٍ.
\end{why}

% ─────────────────────────────────────────────────────────────────────────────
\subsection{فين المشتركون؟ معالِجات الأحداث (\en{EventHandlers})}

السؤال الطبيعيّ: طيب مين بيسمع الحدث ده؟ الإجابة في مكان واحد متوقَّع: تحت
\code{Application/*/EventHandlers/*}. كل مشترِك كلاس بيحقّق \code{INotificationHandler}
لنوع حدث معيّن، وبيحطّ ردّ فعله في دالة \code{Handle}. أمثلة حيّة في النظام:

\begin{itemize}
\item \code{ChatMessageNotificationEventHandler} — بيبعت جرس/بريد لرسالة جديدة.
\item \code{BookingConfirmedEventHandler} — بيبلّغ الطالب إنّ حجزه اتقبل.
\item \code{ApplicationSubmittedNotificationEventHandler} — بيبلّغ إنّ الطلب اتقدّم.
\end{itemize}

ده جسم المشترِك بتاع الشات، مقتصًّا للجوهر:

\begin{english}
\begin{lstlisting}[language={[Sharp]C}]
public sealed class ChatMessageNotificationEventHandler(
    IApplicationDbContext db,
    INotificationDispatcher notifications,
    IChatPresenceQuery chatPresence,
    ILogger<ChatMessageNotificationEventHandler> logger)
    : INotificationHandler<ChatMessageReceivedEvent>          // subscribes to the event
{
    public async Task Handle(ChatMessageReceivedEvent notification, CancellationToken ct)
    {
        // recipient is watching chat now -> the live push already delivered it, skip the bell
        if (chatPresence.IsOnline(notification.RecipientId))
            return;
        // ... look up sender name + build the role-based deep link ...
        await notifications.DispatchAsync(
            notification.RecipientId,
            NotificationType.ChatMessageReceived,
            new NotificationParams { /* CounterpartyName, Preview */ },
            deepLink: deepLink,
            idempotencyKey: $"chat-message:{notification.MessageId:N}",
            ct);
    }
}
\end{lstlisting}
\end{english}

هنشوف إيه سطرًا-بسطر:
\begin{itemize}
\item السطر \code{: INotificationHandler<ChatMessageReceivedEvent>} هو
  الاشتراك نفسه. مجرّد ما الكلاس بيقول «أنا معالِج لهذا الحدث»، \en{MediatR}
  بيسجّله تلقائيًّا وبيناديه لمّا الحدث ينُشَر.
\item دالة \code{Handle} بتاخد الحدث و\code{CancellationToken}، وفيها ردّ الفعل
  كامل. لاحظ إنّ المعالِج ده ما يعرفش مين نشر الحدث ولا ليه — بس بيتصرّف.
\item الحارس \code{if (chatPresence.IsOnline(...)) return;} بيمنع الإزعاج
  المزدوج: لو الطرف فاتح الشات، الرسالة وصلته حيّة، فمفيش داعي لجرس.
\item آخر خطوة: \code{notifications.DispatchAsync(...)} — وهنا المشترِك بيسلّم
  الراية لـ\code{INotificationDispatcher}، اللي شرحنا آليّته بالكامل في فصل
  الإشعارات (منع تكرار \en{$\rightarrow$} نصّ باللغتين \en{$\rightarrow$} قنوات
  \en{$\rightarrow$} حفظ \en{$\rightarrow$} تسليم). المشترِك بيقول «أبلِغ»،
  والموزّع بيتكفّل بالباقي.
\end{itemize}

\begin{note}
لاحظ التغليف بـ\code{try/catch} حوالين \code{DispatchAsync} في الكود الأصليّ:
فشل قناة إشعار ما يصحّش يكسر خطّ الإرسال. لو الإيميل وقع، الرسالة اتخزّنت
خلاص والدفع الحيّ اتبعت — فنسجّل تحذيرًا ونكمّل. ردّ الفعل «إضافيّ»، مش شرطًا
لنجاح الفعل الأساسيّ.
\end{note}

% ─────────────────────────────────────────────────────────────────────────────
\subsection{ليه نعمل «حدث دوميّ» أصلًا؟}

دلوقتي السؤال الصريح: ما دام «احجز» ممكن ببساطة ينادي «ابعت إشعار» على طول، ليه
نلفّ الموضوع بحدث ومشترِك؟ الجواب في كلمة واحدة: الفصل — فصل الاهتمامات
(\en{separation of concerns}). خلّينا نعدّد المكاسب:

\begin{itemize}
\item المعالِج الأساسيّ يفضل نظيفًا: معالِج «احجز» شغلته يعمل الحجز ويحفظه
  وخلاص. ما يعرفش حاجة عن \en{SMTP} ولا \en{SignalR} ولا مقاييس التحليلات.
\item حدث واحد يشغّل عدّة ردود: من نفس \code{BookingConfirmedEvent} ممكن
  يتفرّع إشعار للطالب، وتسجيل مقياس تحليليّ، وبريد تذكير — كلٌّ في مشترِكه.
\item إضافة ردّ جديد بلا لمس المنطق: عايز تضيف «حدّث لوحة التحكّم» بكرة؟
  اكتب مشترِكًا جديدًا لنفس الحدث. مفيش سطر واحد بيتغيّر في معالِج الحجز.
\item اختبار أسهل: تختبر معالِج الحجز إنه نشر الحدث وبس؛ وتختبر كل
  مشترِك لوحده. مفيش تشابك يخلّي اختبارًا واحدًا يلمس نصف النظام.
\end{itemize}

الرسم التالي بيوضّح فكرة \en{fan-out}: نقطة نشر واحدة بتتفرّع لعدّة مشترِكين
مستقلّين. المشترِك الأوّل (الإشعار) موجود فعلًا في النظام، والباقي أمثلة على ردود
ممكن تضيفها لاحقًا بلا ما تلمس معالِج الأمر:

\begin{center}\begin{LTR}
\begin{tikzpicture}[node distance=0.55cm and 1.25cm]
  \node[sbox,text width=2.2cm] (save) {SaveChanges OK};
  \node[gbox,right=1.1cm of save,text width=3.6cm,font=\scriptsize\bfseries] (evt)
    {Publish event\\ChatMessageReceivedEvent};
  \node[srcbox,right=1.4cm of evt,text width=3.0cm] (h1)
    {Notify Handler\\bell + email};
  \node[tbox,above=0.55cm of h1,text width=3.0cm] (h2)
    {Analytics Handler\\usage metric};
  \node[fbox,below=0.55cm of h1,text width=3.0cm] (h3)
    {Digest Handler\\daily summary};
  \draw[arr] (save) -- (evt);
  \draw[arr] (evt.east) -- (h1.west);
  \draw[arr] (evt.east) -- (h2.west);
  \draw[arr] (evt.east) -- (h3.west);
  \node[lbl,below=0.08cm of evt] {one publish};
  \node[lbl,below=0.08cm of h3] {many reactions};
\end{tikzpicture}
\end{LTR}\end{center}

هنشوف إيه في الرسم: العقدة البنفسجيّة في النصّ هي نشر الحدث بعد نجاح الحفظ.
منها بتخرج ثلاثة أسهم لثلاثة مشترِكين — الأخضر (الإشعار) حقيقيّ اليوم، والأزرق
(التحليلات) والبرتقاليّ (الملخّص اليوميّ) أمثلة تُضاف بلا لمس معالِج الأمر. لاحظ:
مصدر واحد، ردود كثيرة، بلا خيط يربط بينهم.

\begin{spotlight}[ركّز هنا لإبهار اللجنة: حدث واحد، ردود كثيرة، بلا تشابك]
لو اللجنة سألت «إزاي بتضيفوا سلوكًا جديدًا من غير ما تكسروا القديم؟» — ده الجواب
النموذجيّ. النظام مبنيّ بحيث تضيف مشترِكًا جديدًا لحدث قائم فيشتغل فورًا، من
غير ما تفتح ملفّ معالِج الأمر الأصليّ أصلًا. ده تطبيق عمليّ لمبدأ «مفتوح للامتداد،
مغلق للتعديل» (\en{Open/Closed}) — وبيتحوّل من شعار في كتاب لخاصّية تلمسها في الكود.
\end{spotlight}

% ─────────────────────────────────────────────────────────────────────────────
\subsection{ربط الدائرة: من الأمر إلى الجرس}

يلّا نوصّل القطع كلّها في مسار واحد كامل. المستخدم يبعت رسالة \en{$\rightarrow$}
معالِج الأمر يحفظ \en{$\rightarrow$} ينشر الحدث \en{$\rightarrow$} المشترِك يلتقطه
\en{$\rightarrow$} المشترِك ينادي \code{DispatchAsync} \en{$\rightarrow$} يتخزّن صفّ
إشعار ويومض الجرس. الرسم ده بيتتبّع الحلقة من أوّلها لآخرها:

\begin{center}\begin{LTR}
\begin{tikzpicture}[node distance=0.7cm and 0.7cm]
  \node[actor,text width=1.7cm] (u) {User};
  \node[pbox,right=0.7cm of u,text width=2.5cm] (cmd) {SendMessage Handler};
  \node[sbox,right=0.6cm of cmd,text width=1.9cm] (save) {SaveChanges};
  \node[gbox,right=0.6cm of save,text width=2.2cm] (pub) {Publish Event};
  \node[srcbox,below=0.85cm of pub,text width=2.6cm] (eh) {EventHandler};
  \node[pbox,left=0.8cm of eh,text width=2.4cm] (disp) {DispatchAsync};
  \node[dbbox,left=0.8cm of disp,text width=2.2cm] (row) {Bell row};
  \draw[arr] (u) -- (cmd);
  \draw[arr] (cmd) -- (save);
  \draw[arr] (save) -- (pub);
  \draw[arr] (pub) -- (eh) node[lbl,right=1pt,midway] {MediatR};
  \draw[arr] (eh) -- (disp);
  \draw[arr] (disp) -- (row);
  \draw[arr] (row.west) to[out=180,in=270] node[lbl,left] {bell blinks} (u.south);
\end{tikzpicture}
\end{LTR}\end{center}

هنشوف إيه: الصفّ العلويّ هو «الفعل» — المستخدم، المعالِج، الحفظ، ثمّ النشر. من عقدة
النشر البنفسجيّة ينزل السهم لـ\en{MediatR} اللي يسلّم الحدث للمشترِك الأخضر
(\code{EventHandler}). المشترِك ينادي \code{DispatchAsync} فيتخزّن صفّ الإشعار
(العقدة ذات الحدّ المزدوج = قاعدة بيانات)، وأخيرًا يومض جرس المستخدم. لاحظ إنّ
النصف السفليّ (الردّ) منفصل تمامًا عن النصف العلويّ (الفعل) — الرابط الوحيد بينهم
هو الحدث.

\begin{qa}[سؤال متوقّع: مش ده تعقيد زيادة؟]
ممكن حدّ يقول: «كان ممكن معالِج الرسالة ينادي الإشعار على طول في سطر، وخلصنا».
صح للحظة، لكن جرّب تضيف تاني وتالت ردّ فعل: كل ردّ هيدخل جوّه معالِج الرسالة، لحدّ
ما يبقى المعالِج «يعرف» عن الإيميل والتحليلات والبريد والدفع — وده بالظبط التشابك
اللي بنهرب منه. الحدث بيدفع تكلفة صغيرة مقدّمًا (نشر + اشتراك) عشان يوفّر تكلفة
ضخمة لاحقًا (كل تعديل بيلمس ملفًّا واحدًا فقط).
\end{qa}

% ─────────────────────────────────────────────────────────────────────────────
\subsection{التدقيق (\en{Audit}): مين عمل إيه وامتى؟}

لو الأحداث بتجاوب على «إيه اللي حصل وإزاي نردّ»، فالتدقيق بيجاوب على سؤال تانيّ
تمامًا: «مين عمل الفعل الحسّاس ده، وامتى، ومن فين؟». استرداد دفعة، تغيير دور
مستخدم، حذف مستند، تعديل إعداد ماليّ — كلّها أفعال محتاجة مساءلة كاملة. ولو
حطّينا كود التسجيل ده جوّه كل معالِج، هيتكرّر وينتشر ويُنسى. الحلّ نفس فلسفة
الأحداث: افصله في طبقة اعتراض واحدة.

الآليّة مكوّنة من قطعتين متعاونتين. الأولى سِمة (\en{attribute}) اسمها
\code{[Auditable]} بتوسم الأمر اللي عايزينه يتسجّل، والتانية سلوك في خطّ أنابيب
\en{MediatR} اسمه \code{AuditBehavior} بيلفّ كل أمر: لو الأمر موسوم، يكتب صفّ تدقيق
بعد نجاح المعالِج؛ ولو مش موسوم، يعدّي بلا عمل. ده جوهر السلوك:

\begin{english}
\begin{lstlisting}[language={[Sharp]C}]
public async Task<TResponse> Handle(
    TRequest request, RequestHandlerDelegate<TResponse> next, CancellationToken ct)
{
    var response = await next(ct);          // run the real handler first

    var attr = (AuditableAttribute?)Attribute.GetCustomAttribute(
        typeof(TRequest), typeof(AuditableAttribute));
    if (attr is null) return response;              // not marked -> no-op

    var targetId = ResolveTargetId(response, request, attr.TargetIdProperty);
    var summary  = BuildSummary(attr.SummaryTemplate, targetId, response, request);
    await audit.WriteAsync(attr.Action, attr.TargetType, targetId,
        beforeJson: null, afterJson: null, summary, ct);   // write the audit row

    return response;
}
\end{lstlisting}
\end{english}

هنشوف إيه:
\begin{itemize}
\item \code{await next(ct)} بينفّذ المعالِج الأصليّ الأوّل. التدقيق بيتكتب
  بعد نجاح الفعل — نفس منطق «سجّل الحقيقة، مش النيّة».
\item \code{Attribute.GetCustomAttribute(...)} بيفحص هل الأمر موسوم بـ
  \code{[Auditable]}. لو \code{null}، السلوك بيرجّع فورًا بلا أيّ أثر — فمفيش
  حمل زائد على الأوامر العاديّة.
\item \code{ResolveTargetId} و\code{BuildSummary} بيستخرجوا «على مين» الفعل ووصفًا
  مقروءًا (زيّ «\en{Refunded payment ...}») من ردّ المعالِج أو الأمر نفسه.
\item \code{audit.WriteAsync(...)} بيكتب الصفّ. لاحظ إنّ الكتابة كلّها ملفوفة بـ
  \code{try/catch} في الأصل: فشل التدقيق ما يصحّش يلغي فعلًا ناجحًا.
\end{itemize}

والوسم نفسه سطر واحد فوق تعريف الأمر. مثال حسّاس — استرداد دفعة:

\begin{english}
\begin{lstlisting}[language={[Sharp]C}]
[Auditable(AuditAction.PaymentRefunded, "Payment",
    TargetIdProperty = nameof(PaymentId),
    SummaryTemplate = "Refunded payment {PaymentId} - amount={AmountCents}")]
public sealed record RefundPaymentCommand(
    Guid PaymentId, long? AmountCents, string? Reason = null) : IRequest<bool>;
\end{lstlisting}
\end{english}

الوسم بيقول لـ\code{AuditBehavior} ثلاث حاجات: نوع الفعل (\code{PaymentRefunded})،
نوع الكيان المستهدَف (\code{Payment})، وأيّ خاصيّة تُؤخَذ كمُعرِّف للهدف
(\code{PaymentId})، وقالب وصف بشريّ. في النظام أكتر من ثمانين أمرًا موسومًا كده —
كل فعل حسّاس تقريبًا: تسجيل الدخول والخروج، تغيير الأدوار، الموافقة والرفض، الحذف،
والاستردادات — بلا سطر تدقيق واحد داخل أيّ معالِج منهم.

طيب، الصفّ اللي بيتكتب بيحتوي إيه بالظبط؟ ده قلب \code{AuditService.WriteAsync}:

\begin{english}
\begin{lstlisting}[language={[Sharp]C}]
var entry = new AuditLog
{
    ActorUserId   = currentUser.UserId,          // WHO   (the actor)
    Action        = action,                      // WHAT  (the verb)
    TargetType    = targetType,                  // on which entity type
    TargetId      = targetId,                     // and which exact row
    BeforeJson    = beforeJson,                   // state BEFORE
    AfterJson     = afterJson,                    // state AFTER
    IpAddress     = Trim(currentUser.IpAddress, 64),      // FROM WHERE
    CorrelationId = Trim(currentUser.CorrelationId, 128), // which request
    OccurredAt    = clock.UtcNow,                 // WHEN
};
db.AuditLogs.Add(entry);   // append-only: only ever Add, never Update or Delete
await db.SaveChangesAsync(ct);
\end{lstlisting}
\end{english}

هنشوف إيه: الصفّ بيجمع كل أبعاد المساءلة في مكان واحد — مَن (\code{ActorUserId})،
ماذا (\code{Action} + \code{TargetType} + \code{TargetId})، متى
(\code{OccurredAt})، من أين (\code{IpAddress} و\code{UserAgent})، وأيّ طلب
ربط الأحداث ببعضها (\code{CorrelationId}). وفيه كمان عمودان \code{BeforeJson} و
\code{AfterJson} لالتقاط لقطة الحالة قبل الفعل وبعده لمّا نحتاج التفصيل ده. الرسم
التالي بيلخّص مسار الاعتراض:

\begin{center}\begin{LTR}
\begin{tikzpicture}[node distance=0.7cm and 0.8cm]
  \node[actor,text width=1.8cm] (adm) {Admin};
  \node[pbox,right=0.75cm of adm,text width=2.6cm] (cmd)
    {[Auditable]\\RefundPayment};
  \node[sbox,right=0.7cm of cmd,text width=2.2cm] (hnd) {Handler runs\\refund};
  \node[gbox,right=0.7cm of hnd,text width=2.4cm] (beh) {AuditBehavior\\after next()};
  \node[dbbox,below=0.9cm of beh,text width=2.3cm] (log) {AuditLog\\append-only};
  \node[sbox,left=1.0cm of log,text width=4.0cm,align=left,font=\scriptsize]
    (fields) {who | what | when\\IpAddress | CorrelationId\\Before / After JSON};
  \draw[arr] (adm) -- (cmd);
  \draw[arr] (cmd) -- (hnd);
  \draw[arr] (hnd) -- (beh);
  \draw[arr] (beh) -- (log) node[lbl,right=1pt,midway] {write};
  \draw[arr] (log) -- (fields) node[lbl,above,midway] {columns};
\end{tikzpicture}
\end{LTR}\end{center}

هنشوف إيه في الرسم: المدير يبعت أمر استرداد موسومًا بـ\code{[Auditable]}. المعالِج
ينفّذ الاسترداد الحقيقيّ، وبعد نجاحه بيلتقط \code{AuditBehavior} الأمر ويكتب صفّ
\code{AuditLog} (الحدّ المزدوج = قاعدة بيانات) بكل الأعمدة على اليسار. لاحظ إنّ
الصفّ \en{append-only}: بنضيفه ومنلمسوش تاني.

\begin{secnote}
ليه \en{append-only} (إضافة فقط، لا تعديل ولا حذف)؟ لأنّ سجلّ التدقيق ما ينفعش
يبقى قابلًا للتزوير — وإلّا مقصده كلّه يسقط. مسار الكود بيضيف صفوفًا بس؛ مفيش أمر
واحد في النظام بيعدّل أو بيمسح صفّ \code{AuditLog}، وجانب القراءة (استعلام السجلّ)
بيقرأ بس. النتيجة: أثر ثابت لا رجعة فيه — حتى المدير اللي عمل الفعل ما يقدرش
يمسح دليله.
\end{secnote}

\begin{spotlight}[ركّز هنا لإبهار اللجنة: كل فعل حسّاس متتبَّع، وبلا رجعة]
اللجنة بتحبّ سؤال «إزاي تعرفوا مين عمل إيه؟». الجواب: أيّ أمر حسّاس بيتوسم بسطر
واحد \code{[Auditable]}، فيلتقطه اعتراض مركزيّ ويكتب صفًّا يحمل الفاعل والفعل
والوقت و\en{IP} ومُعرِّف الطلب، وقابليّة لقطة قبل/بعد — كلّه \en{append-only}. يعني:
مساءلة كاملة، من غير ما نبعثر كود تسجيل في ثمانين معالِجًا، ومن غير ما نثق في إنّ
حدًّا فاكر يكتبه. ده الفرق بين «عندنا \en{logs}» و«عندنا نظام مساءلة».
\end{spotlight}

% ─────────────────────────────────────────────────────────────────────────────
\subsection{الخلاصة}

\begin{summary}[الخلاصة في سطور]
\begin{itemize}
\item الحدث الدوميّ إعلانٌ بأنّ «شيئًا ذا معنى حصل» (\code{BookingConfirmed}،
  \code{ChatMessageReceived}، \code{ApplicationSubmitted})؛ \en{record} ثابت يرث
  \code{DomainEvent} وهو \code{INotification} في \en{MediatR}.
\item يُنشَر بعد \code{SaveChanges} الناجح فقط — عشان مانبلّغش عن تغيير
  اترجّع؛ السطر \code{publisher.Publish(new ...Event(...))} هو نقطة التفرّع.
\item المشترِكون تحت \code{Application/*/EventHandlers/*}؛ كل واحد
  \code{INotificationHandler} لحدث، بيحطّ ردّ فعله في \code{Handle} وبينادي
  \code{DispatchAsync} — بلا معرفة بمَن نشر الحدث.
\item ليه أحداث؟ فصل الاهتمامات: المعالِج الأساسيّ نظيف، حدث واحد يشغّل ردودًا
  كثيرة، إضافة ردّ جديد بلا لمس المنطق، واختبار أبسط.
\item التدقيق: سطر \code{[Auditable]} فوق الأمر + \code{AuditBehavior} يعترض
  بعد النجاح ويكتب \code{AuditLog}: مَن/ماذا/متى/\en{IP}/\en{correlation}، مع
  قابليّة قبل/بعد \en{JSON}، و\en{append-only} — لا تعديل ولا حذف.
\item الفكرة الجامعة: افصل الإعلان عن التصرّف، وافصل المساءلة عن
  المنطق — فيبقى كل شيء مستقلًّا، قابلًا للامتداد، وسهل الاختبار.
\end{itemize}
\end{summary}
